\begin{center}
	\zihao{4}\heiti
	%
	\LARGE{ 中国科学院大学\\$\boldsymbol{2015}$年招收攻读硕士学位研究生入学统一考试试卷\\
	科目名称：电动力学}

\end{center}
\noindent
考生须知：\\
1.本试卷满分为150分,全部考试时间总计180分钟。\\
2.所有答案必须写在答题纸上,写在试题纸上或草稿纸上一律无效。

\vspace*{2ex}
\hrule
\vspace*{4ex}

\fancypagestyle{mypagestyle}{%
	\fancyhf{}% Clear header/footer
	\fancyfoot[OR]{第~\thepage~页，共2页}%
	\fancyfoot[EL]{第~\thepage~页，共2页}%
	\fancyfoot[OL]{科目名称：电动力学 }%
	\fancyfoot[ER]{科目名称：电动力学 }%
	\renewcommand{\headrulewidth}{0pt}% Header rule of .4pt
	\renewcommand{\footrulewidth}{0.5pt}% Header rule of .4pt
}
\fancypagestyle{plain}{%
	\fancyhf{}% Clear header/footer
	\fancyfoot[OR]{第~\thepage~页，共2页}%
	\fancyfoot[EL]{第~\thepage~页，共2页}%
	\fancyfoot[OL]{科目名称：电动力学 }%
	\fancyfoot[ER]{科目名称：电动力学 }%
}

\setcounter{page}{1}%设置页面计数器
\pagestyle{mypagestyle}

\onehalfspacing
% 排版结束，下面是题目

\vspace{0.4cm}

\noindent
一、简答题（共30分，每题10分）
\vspace{0.2cm}

\noindent
1.\hspace{0.8em}写出介质中有源麦克斯韦方程组的微分形式及积分形式。
\vspace{0.3cm}

\noindent
2.\hspace{0.8em}
有一个以理想导体为边界的矩形真空波导，其横截面的两条边长分别为a、b（a>b），写出在该波导中传播的电磁波的波数应满足的条件。对于给定的波模，写出其截止频率的表达式，由此给出该波导管的最低截止频率。\vspace{0.3cm}

\noindent
3.\hspace{0.8em}若在静止坐标系中同一地点两事件的时间间隔为$\Delta \tau$，物体的长度为$l_0$，电荷密度为$\rho_0$，写出在以速度$v$运动的运动坐标系中以上各物理量的变换公式。
\vspace{0.5cm}

\noindent
二、\hspace{0.2em}单项选择题（共30分，每题5分）
\vspace{0.2cm}

\noindent
1.\hspace{0.8em}真空中一个半径为a的不接地、不带电导体球，距球心$R = 2a$处有一点电荷$Q$。该点电荷受到的静电力大小为
\vspace{0.2cm}

A.\hspace{0.8em} $7Q^2 / 180 \pi \varepsilon_0 a^2$ \hspace{10em} B.\hspace{0.8em} $7Q^2 / 288 \pi \varepsilon_0 a^2$\vspace{0.2cm}

C.\hspace{0.8em} $3 Q^2 / 128 \pi \varepsilon_0 a^2$ \hspace{10em} D.\hspace{0.8em} $ 3Q^2 / 16 \pi \varepsilon_0 a^2$
\vspace{0.2cm}

\noindent
2.\hspace{0.8em}无穷长直导线上载有电流I。以该导线为z轴。电流方向为z轴正方向。z>0的区域为真空，z<0的区域充满磁导率为$\mu (\mu > \mu_0 )$的均匀介质。对于界面(z = 0)上的磁化电流，下列说法正确的是
\vspace{0.2cm}

A.\hspace{0.8em}处处为零 \hspace{11em}B.\hspace{0.8em}不等于零，沿角向

C.\hspace{0.8em}不等于零，沿径向向内 \hspace{5em}D.\hspace{0.8em}不等于零，沿径向向外
\vspace{0.2cm}

\noindent
3.\hspace{0.8em}关于位移电流和传导电流，下列说法正确的是\vspace{0.2cm}

A.\hspace{0.8em}两者都消耗电磁能 \hspace{8em}B.\hspace{0.8em}只有前者消耗电磁能

C.\hspace{0.8em}前者与电场有$90^{\circ}$相位差\hspace{5.3em}D.\hspace{0.8em}后者与电场有$90^{\circ}$相位差

\vspace{0.2cm}

\noindent
4.\hspace{0.8em}频率为40Hz的平面电磁波正入射到电导率约为$4 \times 10^8 \mathrm{S} \cdot \mathrm{m}^{-1}$，磁导率约为$4 \pi \times 10^{-7} \mathrm{H} \cdot \mathrm{m}^{-1}$的金属表面，则电磁波的穿透深度约为
\vspace{0.2cm}

A.\hspace{0.8em}0.01cm\hspace{5em} B.\hspace{0.8em}0.4cm\hspace{5em} C.\hspace{0.8em}1.8cm\hspace{5em} D.\hspace{0.8em}5.0cm
\vspace{0.2cm}

\noindent
5.\hspace{0.8em}两个质量和电荷都相同的粒子相向而行发生碰撞，关于该体系下列说法正确的是\vspace{0.2cm}

A.\hspace{0.8em}存在电偶极辐射，但没有磁偶极辐射\vspace{0.1cm}

B.\hspace{0.8em}电偶极辐射与磁偶极辐射都不存在\vspace{0.1cm}

C.\hspace{0.8em}电偶极辐射与磁偶极辐射都存在\vspace{0.1cm}

D.\hspace{0.8em}电偶极辐射与电四极辐射都存在
\vspace{0.2cm}

\noindent
6.\hspace{0.8em}关于远离辐射源$R$处的辐射区电磁场，下列说法正确的是\vspace{0.2cm}

A.\hspace{0.8em}电磁场强度正比于$1/R^2$，能流密度正比于$1/R^4$

B.\hspace{0.8em}电磁场强度正比于$1/R$，能流密度正比于$1/R$

C.\hspace{0.8em}电磁场强度正比于$1/R$，能流密度正比于$1/R^2$

D.\hspace{0.8em}电磁场强度正比于$1/R^2$，能流密度正比于$1/R^2$
\vspace{0.2cm}

\noindent
三、\hspace{0.2em}（共30分）面积为S的平板电容器内充满不均匀介质，极板$A、B$之间的距离为$L$，距离极板$A$为$x$处的介电常数为$\varepsilon(x) = \varepsilon_1 + (\varepsilon_2 - \varepsilon_1) x / L$。在两极板上接上电动势为$V$的电池，即$V_B - V_A = V$（忽略边沿效应），求\vspace{0.2cm}

\noindent
1.距极板$A$为$x$处的电位移矢量和电场强度

\noindent
2.电容器两极板$A$和$B$上的电荷面密度$\sigma_A$和$\sigma_B$

\noindent
3.距极版$A$为$x$处的极化强度和极化电荷密度

\noindent
4.该电容器的电容$C$\vspace{2cm}

\noindent
四、\hspace{0.2em}（共30分）如图所示，有一半径为a，厚度为$\Delta(\Delta \ll a)$的薄电介质圆盘被永久极化，其每一点极化矢量的方向都沿径向向外，极化强度的大小与该点到轴的距离成正比，为$\Vec{P} = p_0 \Vec{r} / 2$，在其中轴线上距圆盘中心z处放一半径为$a(a \ll z)$，载有电流$I$的圆形线圈，圆心在轴线上，线圈的轴线与圆盘的轴线垂直\vspace{0.2cm}

\noindent
1.求此电介质圆盘内的体电荷密度及圆盘边缘的面电荷密度

\noindent
2.若该圆盘以一个恒定的角速度$\omega$绕其轴旋转而同时保持极化强度

\noindent
不变，求此旋转圆盘的磁偶极距 

\noindent
3.求圆盘对线圈的力矩

\vspace{-2.8cm}
\begin{figure}[h!]
    \flushright
    \includegraphics[height=5cm]{test/2015/pic_1.png}
\end{figure}

\noindent
五、\hspace{0.2em}（共30分）如图所示，在z<0区域内充满了电容率为$\varepsilon_1$，磁导率为$\mu_1$的介质，在区域z>0内充满了电容率为$\varepsilon_2$，磁导率为$\mu_2$的介质，有一沿z轴正方向传播的单色平面电磁波，从z<0区域正入射到这两介质的交界面(z = 0)上。已知入射波的电场为$\Large E_i(z,t) = E_{i0}e^{i(k_iz - \omega t)} \Vec{e_x}$
\vspace{0.2cm}

\noindent
1.求界面处反射波和透射波的电场振幅$E_{r0}$，$E_{t0}$与入射波电场振幅$E_{i0}$的比值

\noindent
2.什么条件下反射波与入射波的电场相反？什么条件下无反射波？

\noindent
3.求z<0与z>0区域内电磁波的能流密度对时间的平均值

\vspace{-0.8cm}
\begin{figure}[h!]
    \flushright
    \includegraphics[height=5.5cm]{test/2015/pic_2.png}
\end{figure}

